% ======================================================================
% Chapter 15: Non-Abelian Gauge Invariance
% (book pages 481-503; PDF pages 504-527)
% ======================================================================
\chapter{Non-Abelian Gauge Invariance}

So far in this book we have worked with a rather limited class of quantum
fields and interactions, restricting our attention to scalar field theories,
Yukawa theory, and Quantum Electrodynamics.  It is hardly surprising that
these theories are not sufficient to describe all of the known interactions of
elementary particles.  But what other theories are possible, given that the
Lagrangian of a renormalizable theory can contain no terms of mass dimension
higher than 4?

The most natural theories to try next would be ones with interactions among
vector fields, of the form $A^2(\partial A)$ or $A^4$.  Sensible theories of
this type are difficult to construct, however, because of the negative-norm
states produced by the time component $A^0$ of the vector field operator.  In
Section~5.5 we saw that these negative-norm states cause no difficulty in
QED: They are effectively canceled out by the longitudinal polarization
states, by virtue of the Ward identity.  The Ward identity, in turn, follows
from the invariance of the QED Lagrangian under local gauge transformations.
Perhaps, then, if we can generalize the principle of local gauge invariance,
it will lead us to the construction of other sensible theories of vector
particles.

The goal of this chapter is to do just that.  First we will return briefly to
the study of QED, this time taking the gauge symmetry to be fundamental and
deriving the rest of the theory from this principle.  Then, in Section~15.2,
we will see that the gauge invariance of electrodynamics is only the most
trivial example of an infinite-parameter symmetry, and that the more general
examples lead to other interesting Lagrangians.  These field theories, the
first of which was constructed by Yang and Mills,%
\footnote{C.~N.~Yang and R.~Mills, \emph{Phys.\ Rev.} \textbf{96}, 191 (1954).}
generalize electrodynamics in a profound way.  They are theories of multiple
vector particles, whose interactions are strongly constrained by the symmetry
principle.  In subsequent chapters we will study the quantization of these
theories and their application to the real world of elementary particle
physics.

\section{The Geometry of Gauge Invariance}

In Section~4.1 we wrote down the Lagrangian of Quantum Electrodynamics and
noted the curious fact that it is invariant under a very large group of
transformations \eqref{eq:4.6}, allowing an independent symmetry
transformation at every point in spacetime.  This invariance is the famous
\emph{gauge symmetry} of QED.  From the modern viewpoint, however, gauge
symmetry is not an incidental curiosity, but rather the fundamental principle
that determines the form of the Lagrangian.  Let us now review the elements
of the theory, taking the modern viewpoint.

We begin with the complex-valued Dirac field $\psi(x)$, and stipulate that our
theory should be invariant under the transformation
\begin{equation}
  \psi(x) \to e^{i\alpha(x)}\psi(x) .
  \label{eq:15.1}
\end{equation}
This is a phase rotation through an angle $\alpha(x)$ that varies arbitrarily
from point to point.  How can we write a Lagrangian that is invariant under
this transformation?  As long as we consider terms in the Lagrangian that
have no derivatives, this is easy: We simply write the same terms that are
invariant to global phase rotations.  For example, the fermion mass term
\begin{equation}
  m\bar{\psi}(x)\psi(x)
\end{equation}
is permitted by global phase invariance, and the local invariance gives no
further restriction.

The difficulty arises when we try to write terms including derivatives.  The
derivative of $\psi(x)$ in the direction of the vector $n^\mu$ is defined by the
limiting procedure
\begin{equation}
  n^\mu\partial_\mu\psi(x) = \lim_{\epsilon\to 0}\ \frac{1}{\epsilon}\bigl[
    \psi(x+\epsilon n) - \psi(x) \bigr] .
  \label{eq:15.2}
\end{equation}
However, in a theory with local phase invariance, this definition is not very
sensible, since the two fields that are subtracted, $\psi(x+\epsilon n)$ and
$\psi(x)$, have completely different transformations under the symmetry
\eqref{eq:15.1}.  The quantity $\partial_\mu\psi$, in other words, has no
simple transformation law and no useful geometric interpretation.

In order to subtract the values of $\psi(x)$ at neighboring points in a
meaningful way, we must introduce a factor that compensates for the
difference in phase transformations from one point to the next.  The simplest
way to do this is to define a scalar quantity $U(y,x)$ that depends on the
two points and has the transformation law
\begin{equation}
  U(y,x) \to e^{i\alpha(y)}U(y,x)e^{-i\alpha(x)}
  \label{eq:15.3}
\end{equation}
simultaneously with \eqref{eq:15.1}.  At zero separation, we set
$U(y,y) = 1$; in general, we can require $U(y,x)$ to be a pure phase:
$U(y,x) = \exp[i\phi(y,x)]$.  With this definition, the objects $\psi(y)$ and
$U(y,x)\psi(x)$ have the same transformation law, and we can subtract them in
a manner that is meaningful despite the local symmetry.  Thus we can define a
sensible derivative, called the \emph{covariant derivative}, as follows:
\begin{equation}
  n^\mu D_\mu\psi(x) = \lim_{\epsilon\to 0}\ \frac{1}{\epsilon}\bigl[
    \psi(x+\epsilon n) - U(x+\epsilon n, x)\psi(x) \bigr] .
  \label{eq:15.4}
\end{equation}
To make this definition explicit, we need an expression for the comparator
$U(y,x)$ at infinitesimally separated points.  If the phase of $U(y,x)$ is a
continuous function of the positions $y$ and $x$, then $U(y,x)$ can be
expanded in the separation of the two points:
\begin{equation}
  U(x+\epsilon n, x) = 1 - ie\epsilon n^\mu A_\mu(x) + \mathcal{O}(\epsilon^2) .
  \label{eq:15.5}
\end{equation}
Here we have arbitrarily extracted a constant $e$.  The coefficient of the
displacement $\epsilon n^\mu$ is a new vector field $A_\mu(x)$.  Such a field,
which appears as the infinitesimal limit of a comparator of local symmetry
transformations, is called a \emph{connection}.  The covariant derivative then
takes the form
\begin{equation}
  D_\mu\psi(x) = \partial_\mu\psi(x) + ieA_\mu(x)\psi(x) .
  \label{eq:15.6}
\end{equation}
By inserting \eqref{eq:15.5} into \eqref{eq:15.3}, one finds that $A_\mu$
transforms under this local gauge transformation as
\begin{equation}
  A_\mu(x) \to A_\mu(x) - \frac{1}{e}\partial_\mu\alpha(x) .
  \label{eq:15.7}
\end{equation}
To check that all of these expressions are consistent, we can transform
$D_\mu\psi(x)$ according to Eqs.~\eqref{eq:15.1} and \eqref{eq:15.7}:
\begin{equation}
  D_\mu\psi(x) \to \Bigl( \partial_\mu + ie\Bigl( A_\mu - \frac{1}{e}
    \partial_\mu\alpha \Bigr) \Bigr)\ e^{i\alpha(x)}\psi(x)
  = e^{i\alpha(x)}\bigl( \partial_\mu + ieA_\mu(x) \bigr)\psi(x)
  = e^{i\alpha(x)}D_\mu\psi(x) .
  \label{eq:15.8}
\end{equation}
Thus the covariant derivative transforms in the same way as the field $\psi$
itself, exactly as we constructed it to in the original definition
\eqref{eq:15.4}.

We have now recovered most of the familiar ingredients of the QED
Lagrangian.  From our current viewpoint, however, the definition of the
covariant derivative and the transformation law for the connection $A_\mu$
follow from the postulate of local phase rotation symmetry.  Even the very
existence of the vector field $A_\mu$ is a consequence of local symmetry:
Without it we could not write an invariant Lagrangian involving derivatives
of $\psi$.

More generally, our present analysis gives us a way to construct all possible
Lagrangians that are invariant under the local symmetry.  In any term with
derivatives of $\psi$, replace these with covariant derivatives.  According to
Eq.~\eqref{eq:15.8}, these transform in exactly the same manner as $\psi$
itself.  Therefore any combination of $\psi$ and its covariant derivatives
that is invariant under a global phase rotation (and only these combinations)
will also be locally invariant.

To complete the construction of a locally invariant Lagrangian, we must find
a kinetic energy term for the field $A_\mu$: a locally invariant term that
depends on $A_\mu$ and its derivatives, but not on $\psi$.  This term can be
constructed either integrally, from the comparator $U(y,x)$, or
infinitesimally, from the covariant derivative.

\begin{figure}[t]
  \centering
  \includegraphics[width=0.5\textwidth]{fig15.1.png}
  \caption{Construction of the field strength by comparisons around a small
    square in the (1,2) plane.}
  \label{fig:15.1}
\end{figure}

Working from $U(y,x)$, we will need to extend our explicit formula
\eqref{eq:15.5} to the next term in the expansion in $\epsilon$.  Using the
assumption that $U(y,x)$ is a pure phase and the restriction
$U(x,y)^{\dagger} = U(y,x)$, it follows that
\begin{equation}
  U(x+\epsilon n, x) = \exp\Bigl[ -ie\epsilon n^\mu A_\mu\Bigl( x +
    \tfrac{1}{2}\epsilon n \Bigr) + \mathcal{O}(\epsilon^3) \Bigr] .
  \label{eq:15.9}
\end{equation}
(Relaxing these restrictions introduces additional vector fields into the
theory; this is an unnecessary complication.)  Using this expansion for
$U(y,x)$, we link together comparisons of the phase direction around a small
square in spacetime.  For definiteness, we take this square to lie in the
$(1,2)$-plane, as defined by the unit vectors $\hat{1}$, $\hat{2}$ (see
Fig.~\ref{fig:15.1}).  Define $U(x)$ to be the product of the four
comparisons around the corners of the loop:
\begin{equation}
  U(x) = U(x, x+\epsilon\hat{2})\ U(x+\epsilon\hat{2}, x+\epsilon\hat{1}
    +\epsilon\hat{2})\ U(x+\epsilon\hat{1}+\epsilon\hat{2},
      x+\epsilon\hat{1})\ U(x+\epsilon\hat{1}, x) .
  \label{eq:15.10}
\end{equation}
The transformation law \eqref{eq:15.3} for $U$ implies that $U(x)$ is locally
invariant.  In the limit $\epsilon\to 0$, it will therefore give us a locally
invariant function of $A_\mu$.  To find the form of this function, insert the
expansion \eqref{eq:15.9} to obtain
\begin{equation}
  U(x) = \exp\Bigl\{ -ie\epsilon\Bigl[ -A_2\Bigl( x+\tfrac{\epsilon}{2}\hat{2}
    \Bigr) - A_1\Bigl( x + \epsilon\hat{1} + \tfrac{\epsilon}{2}\hat{2} \Bigr)
      + A_2\Bigl( x + \epsilon\hat{1} + \tfrac{\epsilon}{2}\hat{2} \Bigr)
        + A_1\Bigl( x + \tfrac{\epsilon}{2}\hat{1} \Bigr) \Bigr]
          + \mathcal{O}(\epsilon^4) \Bigr\} .
  \label{eq:15.11}
\end{equation}
When we expand the exponent in powers of $\epsilon$, this reduces to
\begin{equation}
  U(x) = 1 - ie^2\epsilon\bigl[ \partial_1A_2(x) - \partial_2A_1(x) \bigr]
    + \mathcal{O}(\epsilon^3) .
  \label{eq:15.12}
\end{equation}
Therefore the structure
\begin{equation}
  F_{\mu\nu} = \partial_\mu A_\nu - \partial_\nu A_\mu
  \label{eq:15.13}
\end{equation}
is locally invariant.  Of course, $F_{\mu\nu}$ is the familiar electromagnetic
field tensor, and its invariance under \eqref{eq:15.7} can be checked
directly.  The preceding construction, however, shows us the geometrical
origin of the structure of $F_{\mu\nu}$.  Any function that depends on $A_\mu$
only through $F_{\mu\nu}$ and its derivatives is locally invariant.  More
general functions, such as the vector field mass term $m^2A_\mu A^\mu$,
transform under \eqref{eq:15.7} in ways that cannot be compensated and thus
cannot appear in an invariant Lagrangian.

A related argument for the invariance of $F_{\mu\nu}$ can be made using the
covariant derivative.  We have seen above that, if a field has the local
transformation law \eqref{eq:15.1}, then its covariant derivative has the
same transformation law.  Thus the second covariant derivative of $\psi$ also
transforms according to \eqref{eq:15.1}.  The same conclusion holds for the
commutator of covariant derivatives:
\begin{equation}
  [D_\mu, D_\nu]\psi(x) \to e^{i\alpha(x)}[D_\mu, D_\nu]\psi(x) .
  \label{eq:15.14}
\end{equation}
However, the commutator is not itself a derivative at all:
\begin{equation}
  \begin{aligned}
    [D_\mu, D_\nu]\psi &= \bigl[ \partial_\mu\partial_\nu\psi
      + ie(\partial_\mu A_\nu) \psi - ie A_\nu\partial_\mu\psi
      - ie A_\mu\partial_\nu\psi - e^2A_\mu A_\nu\psi \bigr] - (\mu\leftrightarrow\nu)\\
    &= ie\bigl( \partial_\mu A_\nu - \partial_\nu A_\mu \bigr)\cdot\psi .
  \end{aligned}
  \label{eq:15.16}
\end{equation}
That is,
\begin{equation}
  [D_\mu, D_\nu] = ieF_{\mu\nu} .
\end{equation}
On the right-hand side of \eqref{eq:15.14}, the factor $\psi(x)$ accounts for
the entire transformation law, so the multiplicative factor $F_{\mu\nu}$ must
be invariant.  One can visualize the commutator of covariant derivatives as
the comparison of comparisons across a small square; fundamentally,
therefore, this argument is equivalent to that of the previous paragraph.

Whatever the method of proving the invariance of $F_{\mu\nu}$, we have now
assembled all of the ingredients we need to write the most general locally
invariant Lagrangian for the electron field $\psi$ and its associated
connection $A_\mu$.  This Lagrangian must be a function of $\psi$ and its
covariant derivatives, and of $F_{\mu\nu}$ and its derivatives, and must be
invariant to global phase transformations.  Up to operators of dimension 4,
there are only four possible terms:
\begin{equation}
  \mathcal{L}_4 = \bar{\psi}(i\slashed{D})\psi - \frac{1}{4}F_{\mu\nu}F^{\mu\nu}
    - c\epsilon^{\mu\nu\rho\sigma}F_{\mu\nu}F_{\rho\sigma} - m\bar{\psi}\psi .
  \label{eq:15.17}
\end{equation}
By adjusting the normalization of the fields $\psi$ and $A_\mu$, we have set
the coefficients of the first two terms to their standard values.  This
normalization of $A_\mu$ requires the arbitrary scale factor $e$ in our
original definition \eqref{eq:15.5} of $A_\mu$.  The third term violates the
discrete symmetries $P$ and $T$, so we may exclude it if we postulate these
symmetries.%
\footnote{The general systematics of $P$, $C$, and $T$ violation are discussed
in Section~20.3.}  Then $\mathcal{L}_4$ contains only two free parameters,
the scale factor $e$ and the coefficient $m$.

By using operators of dimension 5 and 6, we can form many additional
gauge-invariant combinations:
\begin{equation}
  \mathcal{L}_6 = i c_1\bar{\psi}\sigma^{\mu\nu}F_{\mu\nu}\psi
    + c_2\bigl( \bar{\psi}\psi \bigr)^2 + \cdots .
  \label{eq:15.18}
\end{equation}
More allowed terms appear at each higher order in mass dimension.  But all of
these terms are nonrenormalizable interactions.  In the language of
Section~12.1, they are irrelevant to physics in four dimensions in the limit
where the cutoff is taken to infinity.

\section{The Yang-Mills Lagrangian}

In the previous section, we built up the QED Lagrangian from the postulate of
local phase invariance.  The essential ingredient in this construction was
the covariant derivative, which compensated for the change in the phase of the
fermion field between neighboring points.  In this section, we will generalize
this construction to a broader class of local symmetries.

Consider a theory with a global symmetry under the transformations of a
continuous group.  For definiteness, we take this group to be $SU(2)$---the
group of $2\times2$ unitary matrices with unit determinant.  We represent the
fields of the theory by a pair of complex-valued Dirac fields
\begin{equation}
  \psi(x) = \begin{pmatrix}\psi_1(x)\\ \psi_2(x)\end{pmatrix} .
\end{equation}
We stipulate that our theory should be invariant under the local version of
this symmetry,
\begin{equation}
  \psi(x) \to V(x)\,\psi(x) ,
  \label{eq:15.20}
\end{equation}
where $V(x)$ is a $2\times2$ unitary matrix with unit determinant whose
elements vary arbitrarily with $x$.  As before, the derivative $\partial_\mu\psi$
does not have a simple transformation law, because the fields at neighboring
points transform with different matrices $V$.

To construct a covariant derivative, we again introduce a comparator
$U(y,x)$, which now must be a $2\times2$ matrix with the transformation law
\begin{equation}
  U(y,x) \to V(y)\,U(y,x)\,V^{\dagger}(x) .
  \label{eq:15.22}
\end{equation}
At infinitesimal separation, we expand
\begin{equation}
  U(x+\epsilon n, x) = 1 - ig\epsilon n^\mu A_\mu(x) + \mathcal{O}(\epsilon^2),
  \label{eq:15.23}
\end{equation}
where $A_\mu(x)$ is now a $2\times2$ matrix-valued vector field.  We call such
a quantity a connection.  The covariant derivative is then
\begin{equation}
  D_\mu\psi(x) = \bigl[ \partial_\mu + igA_\mu(x) \bigr]\psi(x) .
  \label{eq:15.24}
\end{equation}
Since $V(x)$ is unitary, the matrices $A_\mu(x)$ can be taken to be Hermitian
and traceless.  A Hermitian, traceless $2\times2$ matrix can be expanded in
the basis of Pauli matrices:
\begin{equation}
  A_\mu(x) = A_\mu^a(x)\ \frac{\sigma^a}{2} .
  \label{eq:15.25}
\end{equation}
Here the index $a$ runs from 1 to 3, and the matrices $\sigma^a/2$ are the
generators of $SU(2)$.  Thus, in the $SU(2)$ gauge theory, there are three
vector fields $A_\mu^a$.

The transformation law of $A_\mu(x)$ follows from the transformation law of
the comparator.  We find
\begin{equation}
  A_\mu(x)\ \frac{\sigma^a}{2} \to V(x)\Bigl( A_\mu(x)\frac{\sigma^a}{2}
    + \frac{1}{g}\partial_\mu \Bigr)\ V^{\dagger}(x) .
  \label{eq:15.27}
\end{equation}
The derivative acts on $V^{\dagger}(x) = \exp(-i\alpha^a\sigma^a/2)$; it is not
so easy to compute this derivative explicitly because the exponent does not
necessarily commute with its derivative.  For infinitesimal transformations
we can expand $V(x)$ to first order in $\alpha$.  In this case we obtain
\begin{equation}
  A_\mu^a\ \frac{\sigma^a}{2} \to \Bigl( A_\mu^a + \frac{1}{g}\partial_\mu\alpha^a
    + \epsilon^{abc}\alpha^b A_\mu^c \Bigr)\ \frac{\sigma^a}{2} + \cdots .
  \label{eq:15.28}
\end{equation}
The last term in this transformation law is new, and arises from the
noncommutativity of the local transformations.  By combining this relation
with the infinitesimal form of the fermion transformation,
\begin{equation}
  \psi \to \Bigl( 1 + i\alpha^a\frac{\sigma^a}{2} \Bigr)\psi + \cdots ,
  \label{eq:15.29}
\end{equation}
we can check the infinitesimal transformation of the covariant derivative:
\begin{equation}
  D_\mu\psi \to \Bigl( 1 + i\alpha^a\frac{\sigma^a}{2} \Bigr)D_\mu\psi ,
  \label{eq:15.30}
\end{equation}
up to terms of order $\alpha^2$.  It is not difficult to check using
\eqref{eq:15.27} and Eq.~(15.21) that, even for finite transformations,
the covariant derivative has the same transformation law as the field on
which it acts.

Using the covariant derivative, we can build the most general gauge-invariant
Lagrangians involving $\psi$.  But to write a complete Lagrangian, we must
also find gauge-invariant terms that depend only on $A_\mu^i$.  To do this, we
construct the analogue of the electromagnetic field tensor.  We will use the
second method of the previous section, working from the commutator of
covariant derivatives.  The transformation law of the covariant derivative
implies that
\begin{equation}
  [D_\mu, D_\nu]\psi(x) \to V(x)[D_\mu, D_\nu]\psi(x) .
  \label{eq:15.31}
\end{equation}
At the same time, by writing out the commutator using formula
\eqref{eq:15.24}, we can show, as in the Abelian case, that $[D_\mu, D_\nu]$
is not a differential operator but merely a multiplicative factor (now a
matrix) acting on $\psi$.  This time, however, there is a new feature: The
last term in the expansion of the commutator no longer vanishes.  Instead, we
find
\begin{equation}
  [D_\mu, D_\nu] = -igF^a_{\mu\nu}\ \frac{\sigma^a}{2} ,
  \label{eq:15.32}
\end{equation}
with
\begin{equation}
  F^a_{\mu\nu} = \partial_\mu A_\nu^a - \partial_\nu A_\mu^a
    + g\epsilon^{abc}A_\mu^b A_\nu^c .
  \label{eq:15.33}
\end{equation}
We can simplify this relation by applying the standard commutation relations
of Pauli matrices:
\begin{equation}
  \Bigl[ \frac{\sigma^a}{2}, \frac{\sigma^b}{2} \Bigr] = i\epsilon^{abc}
    \frac{\sigma^c}{2} .
  \label{eq:15.34}
\end{equation}
Then
\begin{equation}
  F^a_{\mu\nu} = \partial_\mu A_\nu^a - \partial_\nu A_\mu^a
    + g\epsilon^{abc}A_\mu^b A_\nu^c .
  \label{eq:15.35}
\end{equation}
The transformation law for the field strength follows from Eqs.~Eq.~(15.21)
and \eqref{eq:15.31}:
\begin{equation}
  F^a_{\mu\nu}\ \frac{\sigma^a}{2} \to V(x)\ F^b_{\mu\nu}\ \frac{\sigma^b}{2}\
    V^{\dagger}(x) .
  \label{eq:15.36}
\end{equation}
The infinitesimal form is
\begin{equation}
  F^a_{\mu\nu} \to F^a_{\mu\nu} + i\alpha^b F^c_{\mu\nu}\ \epsilon^{bca} .
  \label{eq:15.37}
\end{equation}
Notice that the field strength is no longer a gauge-invariant quantity.  It
cannot be, since there are now three field strengths, each associated with a
given direction of rotation in the abstract space.  However, it is easy to
form gauge-invariant combinations of the field strengths.  For example,
\begin{equation}
  \frac{1}{4}F^a_{\mu\nu}F^{a\mu\nu}
  \label{eq:15.38}
\end{equation}
is a gauge-invariant kinetic energy term for $A_\mu^i$.  Notice that, in
contrast to the case of electrodynamics, this Lagrangian contains cubic and
quartic terms in $A_\mu^i$.  Thus, this Lagrangian describes a nontrivial,
interacting field theory, called \emph{Yang-Mills theory}.  This is the
simplest example of a non-Abelian gauge theory.

To construct a theory of Yang-Mills vector fields interacting with fermions,
we simply add the gauge-field Lagrangian \eqref{eq:15.38} to the familiar
Dirac Lagrangian, with the ordinary derivative of $\psi$ replaced by the
covariant derivative.  The result looks almost identical to the QED
Lagrangian:
\begin{equation}
  \mathcal{L} = \bar{\psi}(i\slashed{D} - m)\psi - \frac{1}{4}(F^a_{\mu\nu})^2 .
  \label{eq:15.39}
\end{equation}
This is the famous Yang-Mills Lagrangian.  Like that of QED, it depends on
two parameters: the scale factor $g$ (which is analogous to the electron
charge) and the fermion mass $m$.  By varying this Lagrangian, we find the
classical equations of motion of the gauge theory.  These are the Dirac
equation for the fermion field and the equation
\begin{equation}
  \partial^\mu F^a_{\mu\nu} + g\epsilon^{abc}A^{b\mu}F^c_{\mu\nu} = gj^a_\nu
  \label{eq:15.40}
\end{equation}
for the vector field.

Everything that we have done for the $SU(2)$ symmetry transformation
\eqref{eq:15.20} generalizes easily to any other continuous group of
symmetries.  The full range of possible symmetry groups is enumerated and
classified in Section~15.4.  For any such group, however, the general
expressions for elements of the Lagrangian are quite similar.  Consider any
continuous group of transformations, represented by a set of $n\times n$
unitary matrices $V$.  Then the basic fields $\psi(x)$ will form an $n$-plet,
and transform according to
\begin{equation}
  \psi(x) \to V(x)\psi(x) ,
  \label{eq:15.41}
\end{equation}
where the $x$ dependence of $V$ makes the transformation local.  In
infinitesimal form, $V(x)$ can be expanded in terms of a set of basic
generators of the symmetry group, which can be represented as Hermitian
matrices $t^a$:
\begin{equation}
  V(x) = 1 + i\alpha^a(x)\,t^a + \mathcal{O}(\alpha^2) .
  \label{eq:15.42}
\end{equation}
Now one can carry through the whole analysis from Eq.~\eqref{eq:15.22} to
Eq.~\eqref{eq:15.33} for a general local symmetry group simply by replacing
\begin{equation}
  \frac{\sigma^a}{2} \to t^a
  \label{eq:15.43}
\end{equation}
at each step of the analysis.

To generalize the explicit expression \eqref{eq:15.35} for the field tensor,
we need to know the commutation relations of the matrices $t^a$.  It is
conventional to write these in the standard form
\begin{equation}
  [t^a, t^b] = if^{abc}t^c ,
  \label{eq:15.44}
\end{equation}
where $f^{abc}$ is a set of numbers called \emph{structure constants}.  This
object replaces $\epsilon^{ijk}$ in Eq.~\eqref{eq:15.34}.  It is conventional
to choose a basis for the matrices $t^a$ such that $f^{abc}$ is completely
antisymmetric; we will prove that this is always possible in Section~15.4.

We can now recapitulate all of our results as follows.  The covariant
derivative associated with the general transformation \eqref{eq:15.41} is
\begin{equation}
  D_\mu = \partial_\mu - igA_\mu^a t^a ;
  \label{eq:15.45}
\end{equation}
it contains one vector field for each independent generator of the local
symmetry.  The infinitesimal transformation laws for $\psi$ and $A_\mu^a$ are
\begin{equation}
  \psi \to \bigl( 1 + i\alpha^a t^a \bigr)\psi;\qquad
  A_\mu^a \to A_\mu^a + \frac{1}{g}\partial_\mu\alpha^a
    + f^{abc}\alpha^b A_\mu^c .
  \label{eq:15.46}
\end{equation}
The finite transformation of $A_\mu^a$ has exactly the form of \eqref{eq:15.27}:
\begin{equation}
  A_\mu^a(x)t^a \to V(x)\Bigl( A_\mu^a(x)t^a + \frac{1}{g}\partial_\mu
    \Bigr)\ V^{\dagger}(x) .
  \label{eq:15.47}
\end{equation}
These transformation laws imply that the covariant derivative of $\psi$ has
the same transformation law as $\psi$ itself.  The field tensor is defined by
\begin{equation}
  [D_\mu, D_\nu] = -igF^a_{\mu\nu}t^a ,
  \label{eq:15.48}
\end{equation}
or more explicitly,
\begin{equation}
  F^a_{\mu\nu} = \partial_\mu A_\nu^a - \partial_\nu A_\mu^a
    + gf^{abc}A_\mu^b A_\nu^c .
  \label{eq:15.49}
\end{equation}
This quantity has the infinitesimal transformation
\begin{equation}
  F^a_{\mu\nu} \to F^a_{\mu\nu} + f^{abc}\alpha^b F^c_{\mu\nu} .
  \label{eq:15.50}
\end{equation}
From Eqs.~\eqref{eq:15.46} and \eqref{eq:15.50}, one can show that any
globally symmetric function of $\psi$, $D_\mu\psi$, and their covariant
derivatives is also locally symmetric, and is therefore a candidate for a
term in a gauge-invariant Lagrangian.  However, there are very few
permissible terms up to dimension 4.  The most general gauge-invariant
Lagrangian that is renormalizable and conserves $P$ and $T$ is again given by
Eq.~\eqref{eq:15.39}.  The corresponding classical equation of motion is
\begin{equation}
  \partial^\mu F^a_{\mu\nu} + gf^{abc}A^{b\mu}F^c_{\mu\nu} = -gj^a_\nu ,
  \label{eq:15.51}
\end{equation}
where
\begin{equation}
  j^{a\mu} = \bar{\psi}\gamma^\mu t^a\psi
  \label{eq:15.52}
\end{equation}
is the global symmetry current of the fermion field.

Notice that the nonlinear terms in the Yang-Mills Lagrangian \eqref{eq:15.39}
appear in the covariant derivative, where they are proportional to $t^a$, and
in the field tensor, where they are proportional to $f^{abc}$.  Thus the form
of the interactions in a non-Abelian gauge theory is dictated by the local
symmetry.  The nonlinear interactions of the vector field with itself are
proportional to the commutators of symmetry generators and thus explicitly
require the non-Abelian nature of the symmetry group.

\section{The Gauge-Invariant Wilson Loop}

In both of the previous sections we made use of the comparator, $U(y,x)$,
which converts the fermion gauge transformation law at point $x$ to that at
point $y$.  So far, in writing expressions for this object, it has sufficed to
assume that $x$ and $y$ are infinitesimally separated.  However, it is
worthwhile to think further about the comparator in the case where $x$ and
$y$ are far apart.  This discussion will give us further insights into the
geometry of gauge invariance, and will reveal some additional useful functions
of the gauge field which we will put to work in Chapter~19.

We first return to the Abelian theory and expand upon our discussion of
$U(y,x)$ in that context.  In Eq.~\eqref{eq:15.10} we constructed a product
of comparators on a path that wound around a small square.  We showed that
this product $U(x)$ is not trivial, even though we eventually return to the
starting point; rather, we found that $U(x)$ differs from 1 by a term
proportional to the electromagnetic field strength and to the area of the
square.  This is a particular case of a general conclusion: The comparator
between two points $x$ and $y$ at finite separation depends on the path taken
from $x$ to $y$.

To explain this statement, it is useful to reverse some of the logic of
Section~15.1.  We begin from the connection $A_\mu$, which we assume to have
the transformation law \eqref{eq:15.7}, and construct $U(z,y)$ as a function
of $A_\mu$ that transforms according to \eqref{eq:15.3}.  It is not difficult
to verify that the expression
\begin{equation}
  U_P(z,y) = \exp\Bigl[ -ie\int_y^z dx^\mu\ A_\mu(x) \Bigr]
  \label{eq:15.53}
\end{equation}
meets this criterion if the integral is taken along any path $P$ that runs
from $y$ to $z$.  This object $U_P(z,y)$ is called the \emph{Wilson line}.%
\footnote{This path-dependent phase was used long before Wilson's work, in
Schwinger's early papers on QED, and in Y.~Aharonov and D.~B\"ohm,
\emph{Phys.\ Rev.} \textbf{115}, 485 (1959).}  Expression \eqref{eq:15.53}
gives an explicit realization of the abstract comparator $U(z,y)$ for points
at finite separation.

A crucial property of the Wilson line is that it depends on the path $P$.  If
$P$ is a closed path that returns to $y$, we obtain the \emph{Wilson loop},
\begin{equation}
  U_P(y,y) = \exp\Bigl[ -ie\oint_P dx^\mu\ A_\mu(x) \Bigr] .
  \label{eq:15.54}
\end{equation}
This quantity is a nontrivial function of $A_\mu$ that is, by construction,
locally gauge invariant.  In fact, all gauge-invariant functions of $A_\mu$
can be thought of as combinations of Wilson loops for various choices of the
path $P$.  To motivate this claim, we use Stokes's theorem to rewrite the
Wilson loop as
\begin{equation}
  \exp\Bigl[ -ie\oint_P dx^\mu\ A_\mu(x) \Bigr]
  = \exp\Bigl[ -\frac{ie}{2}\int_\Sigma d\sigma^{\mu\nu}\ F_{\mu\nu} \Bigr],
  \label{eq:15.55}
\end{equation}
where $\Sigma$ is a surface that spans the closed loop $P$,
$d\sigma^{\mu\nu}$ is an area element on this surface, and $F$ is the field
tensor \eqref{eq:15.13}.  This relation between the Wilson loop and the field
strength is illustrated in Fig.~\ref{fig:15.2}.  Since the Wilson loop is
gauge invariant, this argument gives one more way to visualize the gauge
invariance of the field strength.  Conversely, since (almost) all
gauge-invariant functions of $A_\mu$ can be built up from $F_{\mu\nu}$, this
expression gives weight to the statement that $U_P(y,y)$ is the most general
gauge invariant.

\begin{figure}[t]
  \centering
  \includegraphics[width=0.5\textwidth]{fig15.2.png}
  \caption{The relation between the Wilson loop and the field strength.}
  \label{fig:15.2}
\end{figure}

Both the Wilson line and the Wilson loop can be generalized to the
non-Abelian case.  Here, however, additional subtleties arise when we consider
exponentials of noncommuting matrices.  Let us first construct the Wilson
line, which now transforms according to Eq.~\eqref{eq:15.22}.  It is not
correct to make a straightforward rewriting of \eqref{eq:15.53} with the
integral of $A_\mu^a t^a$ in the exponent, since these matrices do not
necessarily commute at different points.  Instead, we must order these
matrices in a particular way.  We will now give the correct ordering
prescription and then prove its transformation law.

Let $s$ be a parameter of the path $P$, running from 0 at $x=y$ to 1 at
$x=z$.  Then define the Wilson line as the power-series expansion of the
exponential, with the matrices in each term ordered so that higher values of
$s$ stand to the left:
\begin{equation}
  U_P(z,y) = 1 + (-ig)\int_0^1 ds\ \frac{dx^\mu}{ds}\ A_\mu(x(s))
    + (-ig)^2\int_0^1 ds_1\int_0^{s_1} ds_2\ \frac{dx^\mu}{ds_1}\ A_\mu(x(s_1))
      \ \frac{dx^\nu}{ds_2}\ A_\nu(x(s_2)) + \cdots .
  \label{eq:15.56}
\end{equation}
This object is called a \emph{path-ordered exponential}, and is conventionally
written
\begin{equation}
  U_P(z,y) = P\ \exp\Bigl[ -ig\int_y^z dx^\mu\ A_\mu(x) \Bigr] ,
  \label{eq:15.57}
\end{equation}
where the $P$ (for ``path'') instructs us to order the matrices along the
path.  With this definition, one can verify that $U_P(z,y)$ transforms
according to \eqref{eq:15.22}.  The non-Abelian Wilson loop is then
\begin{equation}
  U_P(y,y) = P\ \exp\Bigl[ -ig\oint_P dx^\mu\ A_\mu(x) \Bigr] .
  \label{eq:15.58}
\end{equation}
The trace of this object,
\begin{equation}
  \mathrm{tr}\Bigl[ P\ \exp\Bigl( -ig\oint dx^\mu\ A_\mu \Bigr) \Bigr],
\end{equation}
is a gauge-invariant function of the gauge field, since the trace is invariant
under cyclic permutations.  This quantity plays a central role in the
nonperturbative study of non-Abelian gauge theories.

\section{Basic Facts about Lie Algebras}

At the end of Section~15.2 we saw that the class of non-Abelian gauge
theories is very large.  To work with these theories most efficiently, it is
worthwhile to pause and consider the general properties of the continuous
groups on which they are based.  In this section we will enumerate all the
possible groups that can be used to construct non-Abelian gauge theories.  We
will then compute some numerical factors, built out of group transformation
matrices, that are needed in performing explicit calculations in quantized
gauge theories.%
\footnote{In this section we will state, without proof, some general results
from the theory of continuous groups.  There are many excellent books that
review these mathematical results systematically.  Among these, we recommend
especially Cahn (1984), for a brief but incisive discussion, and
S.~Helgason, \emph{Differential Geometry, Lie Groups, and Symmetric Spaces}
(Academic Press, 1978), which gives an elegant and rigorous account.
R.~Slansky, \emph{Phys.\ Repts.} \textbf{79}, 1 (1981), has compiled an
especially useful set of tables of group-theoretic identities relevant to the
construction of non-Abelian gauge theories.}

To a mathematician, a group is made up of abstract entities that obey certain
algebraic rules.  In quantum mechanics, however, we are interested
specifically in groups of unitary operators that act on the vector space of
quantum states.  We focus our attention on continuously generated groups,
that is, groups that contain elements arbitrarily close to the identity, such
that the general element can be reached by the repeated action of these
infinitesimal elements.  Then any infinitesimal group element $g$ can be
written
\begin{equation}
  g(\alpha) = 1 + i\alpha^a T^a + \mathcal{O}(\alpha^2) .
  \label{eq:15.67}
\end{equation}
The coefficients of the infinitesimal group parameters $\alpha^a$ are
Hermitian operators $T^a$, called the \emph{generators} of the symmetry
group.  A continuous group with this structure is called a \emph{Lie group}.

The set of generators $T^a$ must span the space of infinitesimal group
transformations, so the commutator of generators $T^a$ must be a linear
combination of generators.  Thus the commutation relations of the operators
$T^a$ can be written
\begin{equation}
  [T^a, T^b] = if^{abc}T^c ;
  \label{eq:15.68}
\end{equation}
the numbers $f^{abc}$ are called \emph{structure constants}.  The vector
space spanned by the generators, with the additional operation of
commutation, is called a \emph{Lie algebra}.

The commutation relations \eqref{eq:15.68} and the identity
\begin{equation}
  [T^a, [T^b, T^c]] + [T^b, [T^c, T^a]] + [T^c, [T^a, T^b]] = 0 ,
  \label{eq:15.70}
\end{equation}
called the \emph{Jacobi identity}, define the Lie algebra.  From the
mathematician's viewpoint (considering the generators to be abstract entities
rather than Hermitian operators), the Jacobi identity is an axiom that must
be satisfied in order for a given set of commutation rules to define a Lie
algebra.

The commutation relations of the Lie algebra completely determine the group
multiplication law of an associated Lie group sufficiently close to the
identity.  For large enough transformations, additional global questions come
into play; to give a familiar example, $SU(2)$ and $O(3)$ have the same
commutation relations but different global structure.  However, the
Lagrangian of a non-Abelian gauge theory depends only on the Lie algebra of
the local symmetry group, so we will ignore these global questions from here
on.

\subsection{Classification of Lie Algebras}

For the application to gauge theories, the local symmetry is normally a
unitary transformation of a set of fields.  Thus we are primarily interested
in Lie algebras that have finite-dimensional Hermitian representations, leading
to finite-dimensional unitary representations of the corresponding Lie group.
We will also assume that the number of generators is finite.  Such Lie
algebras are called \emph{compact}, because these conditions imply that the
Lie group is a finite-dimensional compact manifold.

If one of the generators $T^a$ commutes with all of the others, it generates
an independent continuous Abelian group.  Such a group, which has the
structure of the group of phase rotations
\begin{equation}
  \psi \to e^{i\alpha}\psi ,
  \label{eq:15.71}
\end{equation}
we call $U(1)$.  If the algebra contains no such commuting elements, so that
the group contains no $U(1)$ factors, then we call the algebra
\emph{semi-simple}.  If, in addition, the Lie algebra cannot be divided into
two mutually commuting sets of generators, the algebra is \emph{simple}.  A
general Lie algebra is the direct sum of non-Abelian simple components and
additional Abelian generators.

Surprisingly, the basic conditions that a Lie algebra be compact and simple
turn out to be extremely restrictive.  In one of the triumphs of
nineteenth-century mathematics, Killing and Cartan classified all possible
compact simple Lie algebras.  Almost all of these algebras belong to one of
three infinite families, with only five exceptions.  The three infinite
families are the algebras corresponding to the so-called \emph{classical
groups}, whose structures are conveniently defined in terms of particular
matrix representations.  The definitions of the three families of classical
groups are as follows:
\begin{enumerate}
  \item \emph{Unitary transformations of $N$-dimensional vectors.}  Let
    $\xi$ and $\eta$ be complex $N$-vectors.  A general linear transformation
    then has the form
    \begin{equation}
      \eta_a \to U_{ab}\eta_b,\qquad \xi_a \to U_{ab}\xi_b .
      \label{eq:15.72}
    \end{equation}
    We say that this transformation is unitary if it preserves the inner
    product $\eta^{*}_a\xi_a$.  The pure phase transformations
    \begin{equation}
      \xi_a \to e^{i\alpha}\xi_a
      \label{eq:15.73}
    \end{equation}
    form a $U(1)$ subgroup which commutes with all other unitary
    transformations; we remove this subgroup to form a simple Lie group,
    called $SU(N)$; it consists of all $N\times N$ unitary transformations
    satisfying $\det(U) = 1$.  The generators of $SU(N)$ are represented by
    $N\times N$ Hermitian matrices $t^a$, subject to the condition that they
    be orthogonal to the generator of \eqref{eq:15.73}:
    \begin{equation}
      \mathrm{tr}[t^a] = 0 .
      \label{eq:15.74}
    \end{equation}
    There are $N^2-1$ independent matrices satisfying these conditions.
  \item \emph{Orthogonal transformations of $N$-dimensional vectors.}  This
    is the subgroup of unitary $N\times N$ transformations that preserves
    the symmetric inner product
    \begin{equation}
      \eta_a E_{ab}\xi_b,
      \qquad\text{with } E_{ab} = \delta_{ab} .
      \label{eq:15.75}
    \end{equation}
    This is the usual vector product, and so this group is the rotation group
    in $N$ dimensions, $SO(N)$.  (Adding the reflection gives the group
    $O(N)$.)  There is an independent rotation corresponding to each plane in
    $N$ dimensions, so $SO(N)$ has $N(N-1)/2$ generators.
  \item \emph{Symplectic transformations of $N$-dimensional vectors.}  This
    is the subgroup of unitary $N\times N$ transformations, for $N$ even,
    that preserves the antisymmetric inner product
    \begin{equation}
      \eta_a E_{ab}\xi_b,
      \qquad\text{with } E_{ab} = \begin{pmatrix}0&1\\-1&0\end{pmatrix}
      \text{ (in } N/2\times N/2 \text{ blocks)},
      \label{eq:15.76}
    \end{equation}
    where the elements of the matrix are $N/2\times N/2$ blocks.  This group
    is called $Sp(N)$; it has $N(N+1)/2$ generators.
\end{enumerate}

Beyond these three families, there are five more \emph{exceptional} Lie
algebras, denoted in Cartan's classification system as $G_2$, $F_4$, $E_6$,
$E_7$, and $E_8$.  Of these, $E_6$ and $E_8$ have been applied as local
symmetry groups in interesting unified models of the fundamental
interactions.  However, we will not consider these exceptional groups further
in this book.  In fact, most of our examples will involve only $SU(N)$
groups.

\subsection{Representations}

Once we have specified the local symmetry group, the fields that appear in
the Lagrangian most naturally transform according to a finite-dimensional
unitary representation of this group.  Thus we might next ask how to
systematically find all such representations of any given Lie group.  Recall
that for the group $SU(2)$, the representations can be constructed directly
from the commutation relations, using the raising and lowering operators
$J^+$ and $J^-$.  This construction can be generalized to find the
finite-dimensional representations of any compact Lie algebra.  In this book,
however, we will work with relatively simple representations whose structure
we can work out by less formal methods.

Before discussing representations of Lie algebras, we should review some
general aspects of group representations.  Given a symmetry group $G$, a
finite-dimensional unitary representation of the group's Lie algebra is a set
of $d\times d$ Hermitian matrices $t^a$ that satisfy the commutation
relations \eqref{eq:15.68}.  The size $d$ is the \emph{dimension} of the
representation.  An arbitrary representation can generally be decomposed by
finding a basis in which all representation matrices are simultaneously
block-diagonal.  Through this change of basis, we can write the
representation as the direct sum of \emph{irreducible representations}.  We
denote the representation matrices in the irreducible representation $r$ by
$t^a_r$.

It is standard practice to adopt a normalization convention for the matrices
$t^a$ based on traces of their products.  If the Lie algebra is semi-simple,
the matrices $t^a_r$ themselves are traceless.  Consider, however, the trace
of the product of two generator matrices:
\begin{equation}
  \mathrm{tr}\bigl[ t^a_r t^b_r \bigr] = D^{ab} .
  \label{eq:15.77}
\end{equation}
As long as the generator matrices are Hermitian, the matrix $D^{ab}$ is
positive definite.  Let us choose a basis for the generators $T^a$ so that
this matrix is proportional to the identity.  It can be shown that, once this
is done for one irreducible representation, it is true for all irreducible
representations.  Thus, in this basis,
\begin{equation}
  \mathrm{tr}\bigl[ t^a_r t^b_r \bigr] = C(r)\,\delta^{ab} ,
  \label{eq:15.78}
\end{equation}
where $C(r)$ is a constant for each representation $r$.  Equation
\eqref{eq:15.78} and the commutation relations \eqref{eq:15.68} yield the
following representation of the structure constants:
\begin{equation}
  f^{abc} = -\frac{i}{C(r)}\,\mathrm{tr}\bigl( \bigl[ t^a_r, t^b_r \bigr]
    t^c_r \bigr) .
  \label{eq:15.79}
\end{equation}
This equation implies that $f^{abc}$ is totally antisymmetric.

For each irreducible representation $r$ of $G$, there is an associated
\emph{conjugate representation} $\bar{r}$.  The representation $r$ yields the
infinitesimal transformation
\begin{equation}
  \phi \to \bigl( 1 + i\alpha^a t^a_r \bigr)\phi .
  \label{eq:15.80}
\end{equation}
The complex conjugate of this transformation,
\begin{equation}
  \phi^{*} \to \bigl( 1 - i\alpha^a (t^a_r)^{*} \bigr)\phi^{*} ,
  \label{eq:15.81}
\end{equation}
must also be the infinitesimal element of a representation of $G$.  Thus the
conjugate representation to $r$ has representation matrices
\begin{equation}
  t^a_{\bar{r}} = -(t^a_r)^{*} = -(t^a_r)^T .
  \label{eq:15.82}
\end{equation}
Since $\phi^{*}\phi$ is invariant to unitary transformations, it is possible
to combine fields transforming in the representations $r$ and $\bar{r}$ to
form a group invariant.

It is possible that the representation $\bar{r}$ may be equivalent to $r$, if
there is a unitary transformation $U$ such that
$t^a_{\bar{r}} = Ut^a_rU^{\dagger}$.  If so, the representation $r$ is
\emph{real}.  In this case, there is a matrix $G_{ab}$ such that, if $\eta$ and
$\xi$ belong to the representation $r$, the combination
$G_{ab}\eta_a\xi_b$ is an invariant.  It is sometimes useful to distinguish
the case in which $G_{ab}$ is symmetric from that in which $G_{ab}$ is
antisymmetric.  In the former case the representation is \emph{strictly
real}; in the latter case it is \emph{pseudoreal}.  Both cases occur already
in $SU(2)$: The invariant combination of two vectors is $v^aw^a$, so the
vector is a real representation; the invariant combination of two spinors is
$\epsilon^{ab}\chi_a\psi_b$, so the spinor is a pseudoreal representation.

With this language we can discuss the simplest representations of the
classical groups.  In $SU(N)$, the basic irreducible representation (often
called the \emph{fundamental representation}) is the $N$-dimensional complex
vector.  For $N > 2$ this representation is complex, so that there is a
second, inequivalent, representation $\bar{N}$.  (In $SU(2)$ this
representation is the pseudoreal spinor representation.)  In $SO(N)$, the
basic $N$-dimensional vector is a (strictly) real representation.  In
$Sp(N)$, the $N$-dimensional vector is a pseudoreal representation.

Another irreducible representation, present for any simple Lie algebra, is
the one to which the generators of the algebra belong.  This representation
is called the \emph{adjoint representation} and denoted by $r = G$.  The
representation matrices are given by the structure constants:
\begin{equation}
  (t^b_G)_{ac} = if^{abc} .
  \label{eq:15.83}
\end{equation}
With this definition, the statement that $t^b_G$ satisfies the Lie algebra
\begin{equation}
  \bigl( [t^b_G, t^c_G] \bigr)_{ae} = if^{bcd}(t^d_G)_{ae}
  \label{eq:15.84}
\end{equation}
is just a rewriting of the Jacobi identity \eqref{eq:15.70}.  Since the
structure constants are real and antisymmetric,
$t^a_G = -(t^a_G)^{*}$; thus the adjoint representation is always a real
representation.  From the descriptions of the Lie groups given above, the
dimension of the adjoint representation $d(G)$ is given, for the classical
groups, by
\begin{equation}
  d(G) = \begin{cases}
    N^2-1 & \text{for } SU(N),\\
    N(N-1)/2 & \text{for } SO(N),\\
    N(N+1)/2 & \text{for } Sp(N) .
  \end{cases}
  \label{eq:15.85}
\end{equation}

The identification of $f^{abc}$ as a representation matrix allows us to gain
further insight into some of the quantities introduced in Section~15.2.  The
covariant derivative acting on a field in the adjoint representation is
\begin{equation}
  \begin{aligned}
    (D_\mu\phi)^a &= \partial_\mu\phi^a - igA_\mu^b(t^b_G)_{ac}\phi^c\\
    &= \partial_\mu\phi^a + gf^{abc}A_\mu^b\phi^c .
  \end{aligned}
  \label{eq:15.86}
\end{equation}
Thus we can recognize the infinitesimal form of the gauge transformation of
the vector field in \eqref{eq:15.46} as the motion of the field under the
adjoint representation.

Another important group-theoretic quantity is the \emph{quadratic Casimir
operator}.  For a given irreducible representation $r$, we define
\begin{equation}
  C_2(r) = \sum_a t^a_r\,t^a_r .
  \label{eq:15.92}
\end{equation}
By Schur's lemma, this matrix commutes with all of the generators $t^a_r$, and
so, for an irreducible representation, it must be proportional to the
identity:
\begin{equation}
  C_2(r) = C_2(r)\cdot\mathbb{1} .
  \label{eq:15.93}
\end{equation}
The proportionality constant $C_2(r)$ is the eigenvalue of the Casimir
operator in the representation $r$.  Taking the trace of \eqref{eq:15.92}
and using \eqref{eq:15.78}, we find
\begin{equation}
  C_2(r)\cdot d(r) = C(r)\cdot d(G) ,
  \label{eq:15.94}
\end{equation}
where $d(r)$ is the dimension of the representation $r$.  For the adjoint
representation, the Casimir operator is
\begin{equation}
  C_2(G) = C(G) .
\end{equation}
For the fundamental representation of $SU(N)$, one can compute $C_2(r)$
directly.  The fundamental representation of $SU(N)$ acts on $N$-dimensional
complex vectors, and its generators transform as a general traceless
$N\times N$ tensor; the matrices that effect these transformations make up
the adjoint representation of $SU(N)$.  In this case, the decomposition of
the product of the fundamental and conjugate representations gives
\begin{equation}
  N\times\bar{N} = 1 + (N^2-1) .
  \label{eq:15.102}
\end{equation}
For this decomposition, Eqs.~Eq.~(15.100) and Eq.~(15.101) imply
the identity
\begin{equation}
  \Bigl( 2C_2(r)\cdot N \Bigr) \cdot N^2 = 0 + C_2(G)\cdot(N^2-1) .
  \label{eq:15.103}
\end{equation}
Thus, for $SU(N)$,
\begin{equation}
  C_2(G) = C(G) = N .
  \label{eq:15.104}
\end{equation}
Some additional examples of the computation of quadratic Casimir operators
are given in Problem~15.5.  However, the examples we have discussed in this
section, combined with the basic group-theoretic concepts that we have
reviewed, already provide enough material to carry out the most important
computations of physical interest in non-Abelian gauge theories.

\problems

\textbf{15.1} \emph{Brute-force computations in $SU(3)$.}  The standard basis
for the fundamental representation of $SU(3)$ is
\begin{equation}
  \lambda^1 = \begin{pmatrix}0&1&0\\1&0&0\\0&0&0\end{pmatrix},\qquad
  \lambda^2 = \begin{pmatrix}0&-i&0\\i&0&0\\0&0&0\end{pmatrix},\qquad
  \lambda^3 = \begin{pmatrix}1&0&0\\0&-1&0\\0&0&0\end{pmatrix},\qquad
  \cdots
\end{equation}
\begin{enumerate}[label=(\alph*)]
  \item Explain why there are exactly eight matrices in the basis.
  \item Evaluate all the commutators of these matrices, to determine the
    structure constants of $SU(3)$.  Show that, with the normalizations used
    here, $f^{abc}$ is totally antisymmetric.  (This exercise is tedious; you
    may wish to check only a representative sample of the commutators.)
  \item Check the orthogonality condition \eqref{eq:15.78}, and evaluate the
    constant $C(r)$ for this representation.
  \item Compute the quadratic Casimir operator $C_2(r)$ directly from its
    definition \eqref{eq:15.92}, and verify the relation \eqref{eq:15.94}
    between $C_2(r)$ and $C(r)$.
\end{enumerate}

\textbf{15.2} Write down the basis matrices of the adjoint representation of
$SU(2)$.  Compute $C(G)$ and $C_2(G)$ directly from their definitions
\eqref{eq:15.78} and \eqref{eq:15.92}.

\textbf{15.3} \emph{Coulomb potential.}
\begin{enumerate}[label=(\alph*)]
  \item Using functional integration, compute the expectation value of the
    Wilson loop in pure quantum electrodynamics without fermions.  Show that
    \begin{equation}
      \langle U_P(z,z) \rangle = \exp\Bigl[ -\frac{e^2}{2}\oint_P dx^\mu\;
        \oint_P dy^\nu\ g_{\mu\nu}\ \frac{1}{8\pi^2(x-y)^2} \Bigr] ,
    \end{equation}
    with $x$ and $y$ integrated around the closed curve $P$.
  \item Consider the Wilson loop of a rectangular path of (spacelike) width
    $R$ and (timelike) length $T$, $T\gg R$.  Compute the expectation value of
    the Wilson loop in this limit and compare to the general expression for
    time evolution,
    \begin{equation}
      \langle U_P \rangle = \exp[-iE(R)T] ,
    \end{equation}
    where $E(R)$ is the energy of the electromagnetic sources corresponding
    to the Wilson loop.  Show that the potential energy of these sources is
    just the Coulomb potential, $V(R) = -e^2/4\pi R$.
  \item Assuming that the propagator of the non-Abelian gauge field is given
    by the Feynman gauge expression
    \begin{equation}
      \frac{-i\delta^{ab}g_{\mu\nu}}{q^2+i\epsilon} ,
    \end{equation}
    compute the leading-order expectation value of a small rectangular
    Wilson loop in a non-Abelian gauge theory, and show that the result is
    proportional to the Casimir operator $C_2(r)$ of the representation of
    the source.
\end{enumerate}
